\documentclass[11pt]{article}
\usepackage[margin=1.05in]{geometry}
\usepackage{amsmath,amssymb,amsthm,mathrsfs}
\usepackage{hyperref}

\newtheorem{thm}{Theorem}[section]
\newtheorem{lem}[thm]{Lemma}
\newtheorem{prop}[thm]{Proposition}
\newtheorem{cor}[thm]{Corollary}
\theoremstyle{definition}
\newtheorem{defn}[thm]{Definition}
\newtheorem{rem}[thm]{Remark}
\newtheorem{ex}[thm]{Example}

\newcommand{\Z}{\mathbb{Z}}
\newcommand{\Saff}{\widetilde{S}_{n}}
\newcommand{\XX}{\mathcal{X}}

\title{Deleting a bracketed pair preserves length-additivity\\
\large (and the bracket is not needed)}
\author{Clio}
\date{23 September 2026}

\begin{document}
\maketitle

\begin{center}\small
\begin{tabular}{ll}
\textbf{Author:} Clio &
\textbf{Date:} 23 September 2026 \\
\textbf{Recipient:} Robin Langer &
\textbf{Commit:} \texttt{clio-vega/proofs@eb9d081} \\
\end{tabular}\\[2pt]
\textbf{Title:} Deleting a bracketed pair preserves length-additivity (and the bracket is
not needed).\\[1pt]
\emph{This note is new at the commit named above; it depends on that commit only through
\cite{Clio2}, whose Lemmas 2.2 and 2.3 are restated here in full with their proofs, so the
present note is mathematically self-contained.}
\end{center}

\begin{abstract}
Let $S,T\subsetneq\Z/n$ and suppose the factorisation $u_Su_T$ of cyclically decreasing
elements of the affine symmetric group is length-additive. Question Q232 asks whether, when
$i\in S$ and $i+1\in T$, the reduced pair $(S\setminus\{i\},\,T\setminus\{i+1\})$ is again
length-additive. We prove that it is. In fact we prove more: the hypothesis $i+1\in T$ is
never used, so for \emph{every} $i\in S$ the pair
$(S\setminus\{i\},\,T\setminus\{i+1\})$ is length-additive, whether or not $i+1$ lies in $T$.
The proof is a three-case check against the length-additivity criterion, made possible by the
observation that deleting $i$ from $S$ changes $u_S$ at exactly two positions --- lowering the
value at one and raising it at the other --- and that the raised position is precisely the one
deleted from $T$.
\end{abstract}

\section{Conventions, and the two facts we import}\label{sec:conv}

Let $n\ge2$. The affine symmetric group $\Saff$ is the group of bijections $w:\Z\to\Z$ with
$w(j+n)=w(j)+n$ and $\sum_{j=1}^{n}w(j)=\binom{n+1}{2}$; for $i\in\Z/n$, $s_i$ interchanges
$i+jn$ and $i+1+jn$ for all $j$. Products act by $(ab)(j)=a(b(j))$ and
$\ell(ws_i)=\ell(w)+1$ iff $w(i)<w(i+1)$. Every $w\in\Saff$ is $n$-equivariant:
\begin{equation}\label{eq:equiv}
  w(x+n)=w(x)+n\qquad(x\in\Z).
\end{equation}

A subset $S\subseteq\Z/n$ is identified with its $n$-periodic lift $\{j\in\Z: j\bmod n\in S\}$,
written again $S$. A \emph{run} of $S$ is a maximal integer interval $[m,M]\subseteq S$; a
\emph{gap} is an element of $\Z\setminus S$. If $S\subsetneq\Z/n$ then the gap set is nonempty
and $n$-periodic, so gaps recur with period $n$ and
\begin{equation}\label{eq:runlen}
  \text{every run of a proper }S\text{ has length }\le n-1 .
\end{equation}
The \emph{blocks} of $S$ are the intervals $[g+1,g']$ for consecutive gaps $g<g'$; they
partition $\Z$, and a run $[m,M]$ gives the block $[m,M+1]$. For $S\subsetneq\Z/n$ we write
$u_S$ for the unique cyclically decreasing element with support $S$. We call $(S,T)$
\emph{additive} if $S,T\subsetneq\Z/n$ and $\ell(u_Su_T)=|S|+|T|$.

We import two facts from \cite[Lems.~2.2, 2.3]{Clio2}, restated here with their proofs.

\begin{lem}[Block model, {\cite[Lem.~2.2]{Clio2}}]\label{lem:blocks}
Let $S\subsetneq\Z/n$. Then $u_S$ rotates each block downwards by one step:
\begin{equation}\label{eq:rot}
  u_S(g+1)=g',\qquad u_S(j)=j-1\ \ (g+1<j\le g')
\end{equation}
for every block $[g+1,g']$. Equivalently $u_S(j)=j-1$ if $j-1\in S$, and
$u_S(j)=\min\{c\ge j: c\notin S\}$ otherwise. Moreover $\ell(u_S)=|S|$.
\end{lem}

\begin{proof}
For a run $[m,M]$ the product $s_Ms_{M-1}\cdots s_m$ is a reduced word in distinct generators
in which $s_{i+1}$ precedes $s_i$ whenever both occur; as a permutation it sends $m\mapsto M+1$
and $j\mapsto j-1$ for $m<j\le M+1$, i.e.\ it is the rotation \eqref{eq:rot} of the block
$[m,M+1]$ and the identity elsewhere. Distinct runs are separated by a gap, so their blocks are
disjoint and the rotations commute; the product is the stated permutation, of length
$\sum(M-m+1)=|S|$. See \cite[Lem.~2.2]{Clio2} for reducedness and uniqueness, which we do not
use below.
\end{proof}

\begin{lem}[Length-additivity criterion, {\cite[Lem.~2.3]{Clio2}}]\label{lem:add}
Let $w\in\Saff$ and $X\subsetneq\Z/n$. Then $\ell(wu_X)=\ell(w)+|X|$ if and only if for every
run $[m,M]$ of $X$,
\[
  w(M+1)>w(j)\qquad\text{for all }j\in[m,M].
\]
\end{lem}

\begin{proof}
For a single run, $u_{[m,M]}=s_Ms_{M-1}\cdots s_m$. Multiplying on the right one letter at a
time, $\ell(ws_M)=\ell(w)+1$ iff $w(M)<w(M+1)$; the element $ws_M$ has $M+1$ and $M$ swapped in
position, so $\ell(ws_Ms_{M-1})$ increases iff $(ws_M)(M-1)<(ws_M)(M)$, i.e.\ iff
$w(M-1)<w(M+1)$; inductively the $j$-th step increases the length iff $w(M-j+1)<w(M+1)$. Hence
the run contributes $M-m+1$ iff $w(M+1)>w(j)$ for all $j\in[m,M]$. Distinct runs act on
disjoint position-blocks and commute, and each condition involves only $w$, so the conditions
are simultaneous. By \eqref{eq:equiv} the condition for a run is invariant under replacing
$[m,M]$ by $[m+kn,M+kn]$, so it suffices to check one lift of each run.
\end{proof}

\section{Deleting one element from \texorpdfstring{$S$}{S}}\label{sec:delete}

Fix $S\subsetneq\Z/n$ and $i\in S$; choose an integer lift of $i$, again written $i$, and let
$[p,q]$ be the run of $S$ containing $i$, so
\begin{equation}\label{eq:pq}
  p\le i\le q,\qquad p-1\notin S,\qquad q+1\notin S,\qquad q-p\le n-2
\end{equation}
by \eqref{eq:runlen}. Write $S'=S\setminus\{i\}$ throughout this section.

\begin{lem}[Two-position perturbation]\label{lem:perturb}
The maps $u_S$ and $u_{S'}$ agree at every $x\in\Z$ with $x\not\equiv p$ and
$x\not\equiv i+1\pmod n$, and
\[
  u_S(p)=q+1,\quad u_{S'}(p)=i,\qquad u_S(i+1)=i,\quad u_{S'}(i+1)=q+1 .
\]
In particular $u_{S'}(p)<u_S(p)$ and $u_{S'}(i+1)>u_S(i+1)$.
\end{lem}

\begin{proof}
The gap set of $S'$ is that of $S$ together with the residue class of $i$. Hence every block of
$S$ other than $[p,q+1]$ and its translates is a block of $S'$, while $[p,q+1]$ splits into the
two blocks $[p,i]$ and $[i+1,q+1]$ (both nonempty, as $p\le i$ and $i+1\le q+1$). Outside the
translates of $[p,q+1]$ the two elements are therefore given by the same rotations and agree.
Inside, Lemma~\ref{lem:blocks} gives
\[
  u_S:\ p\mapsto q+1,\quad j\mapsto j-1\ (p<j\le q+1),
\]
\[
  u_{S'}:\ p\mapsto i,\quad j\mapsto j-1\ (p<j\le i),\qquad
           i+1\mapsto q+1,\quad j\mapsto j-1\ (i+1<j\le q+1).
\]
These agree at every $j\in[p,q+1]$ except $j=p$ and $j=i+1$, which are distinct because
$p\le i<i+1$. The displayed values are read off directly, and
$u_{S'}(p)=i\le q<q+1=u_S(p)$, $u_{S'}(i+1)=q+1>i=u_S(i+1)$.
\end{proof}

\begin{lem}[The two moved positions are incongruent]\label{lem:incong}
$p\not\equiv i+1 \pmod n$.
\end{lem}

\begin{proof}
By \eqref{eq:pq}, $i-p\le q-p\le n-2$, so $i+1-n\le p-1<p$; and $p\le i<i+1$. Hence
$i+1-n<p<i+1$, an open interval of length $n$ whose only points congruent to $i+1$ modulo $n$
are its endpoints. So $p\not\equiv i+1$.
\end{proof}

\begin{lem}[Nothing below $p$ climbs above $p-1$]\label{lem:below}
For every $x\le p-1$ we have $u_S(x)\le p-1$ and $u_{S'}(x)\le p-1$.
\end{lem}

\begin{proof}
By \eqref{eq:pq}, $p-1$ is a gap of $S$, hence also of $S'\subseteq S$. Let $[g+1,g']$ be the
block of $S$ (resp.\ of $S'$) containing $x$, so $g<x\le g'$ with $g,g'$ consecutive gaps. If
$g'>p-1$ then $g<x\le p-1<g'$, so the gap $p-1$ lies strictly between the consecutive gaps $g$
and $g'$ --- impossible. Hence $g'\le p-1$. By Lemma~\ref{lem:blocks} the block is preserved,
so $u(x)\in[g+1,g']$ and $u(x)\le g'\le p-1$.
\end{proof}

\section{The theorem}\label{sec:thm}

\begin{thm}[Deletion]\label{thm:main}
Let $n\ge2$, let $(S,T)$ be additive, and let $i\in S$. Then
\[
  \ell\bigl(u_{S\setminus\{i\}}\,u_{T\setminus\{i+1\}}\bigr)
  \;=\;\bigl|S\setminus\{i\}\bigr|+\bigl|T\setminus\{i+1\}\bigr| ,
\]
i.e.\ $(S\setminus\{i\},\,T\setminus\{i+1\})$ is additive.
\end{thm}

\begin{proof}
Write $S'=S\setminus\{i\}$ and $T'=T\setminus\{i+1\}$, and fix an integer lift $i$ of $i$ with
its $S$-run $[p,q]$ as in \S\ref{sec:delete}. Both $S'$ and $T'$ are proper subsets of $\Z/n$,
and $\ell(u_{S'})=|S'|$ by Lemma~\ref{lem:blocks}. By Lemma~\ref{lem:add} applied with
$w=u_{S'}$ and $X=T'$, it suffices to prove
\begin{equation}\label{eq:goal}
  u_{S'}(M'+1)>u_{S'}(j)\qquad\text{for every run }[m',M']\text{ of }T'\text{ and every }
  j\in[m',M'].
\end{equation}
If $T'=\emptyset$ there is nothing to prove. Fix a run $[m',M']$ of $T'$; by \eqref{eq:equiv}
we may translate it by any multiple of $n$, since \eqref{eq:goal} is invariant under that
translation. By \eqref{eq:runlen},
\begin{equation}\label{eq:short}
  m'\ \ge\ M'-n+2 .
\end{equation}

Two remarks are used repeatedly. First, every $j\in[m',M']$ lies in $T'$, so
\begin{equation}\label{eq:notiplus1}
  j\not\equiv i+1 \pmod n \qquad (j\in[m',M']),
\end{equation}
because $T'$ is the periodic lift of a set not containing the residue of $i+1$. Second, by
Lemma~\ref{lem:perturb} and \eqref{eq:notiplus1},
\begin{equation}\label{eq:mono}
  u_{S'}(j)\le u_S(j)\qquad (j\in[m',M']),
\end{equation}
with equality unless $j\equiv p$.

We distinguish cases according to the residue of $M'+1$. By Lemma~\ref{lem:incong} the two
conditions $M'+1\equiv p$ and $M'+1\equiv i+1$ are mutually exclusive.

\medskip
\noindent\textbf{Case 2: $M'+1\equiv p\pmod n$.}
Translate the run so that $M'+1=p$, i.e.\ $M'=p-1$. Then $u_{S'}(M'+1)=u_{S'}(p)=i$ by
Lemma~\ref{lem:perturb}. Every $j\in[m',M']$ satisfies $j\le p-1$, so
$u_{S'}(j)\le p-1$ by Lemma~\ref{lem:below}. Since $p\le i$ we get $u_{S'}(j)\le p-1\le i-1<i
=u_{S'}(M'+1)$, which is \eqref{eq:goal}.

\medskip
\noindent In the remaining cases $M'+1\not\equiv p$, so Lemma~\ref{lem:perturb} gives
\begin{equation}\label{eq:up}
  u_{S'}(M'+1)\ \ge\ u_S(M'+1).
\end{equation}
Since $M'+1\notin T'$, either $M'+1\notin T$ or else $M'+1\in T$ and $M'+1\equiv i+1$.

\medskip
\noindent\textbf{Case 1a: $M'+1\notin T$.}
Then $M'$ is the top of the run $[a,b]$ of $T$ containing $[m',M']$ (such a run exists because
$T'\subseteq T$ and $[m',M']$ is an interval in $T'$), i.e.\ $b=M'$. Additivity of $(S,T)$ and
Lemma~\ref{lem:add} applied with $w=u_S$ and $X=T$ give $u_S(b+1)>u_S(j)$ for all $j\in[a,b]$,
hence in particular
\[
  u_S(M'+1)>u_S(j)\qquad (j\in[m',M']).
\]
Combining with \eqref{eq:up} on the left and \eqref{eq:mono} on the right,
$u_{S'}(M'+1)\ge u_S(M'+1)>u_S(j)\ge u_{S'}(j)$, which is \eqref{eq:goal}.

\medskip
\noindent\textbf{Case 1b: $M'+1\in T$ and $M'+1\equiv i+1\pmod n$.}
Translate the run so that $M'+1=i+1$, i.e.\ $M'=i$. Then $u_{S'}(M'+1)=u_{S'}(i+1)=q+1$ by
Lemma~\ref{lem:perturb}, and by \eqref{eq:short} we have $m'\ge i-n+2$. Let $j\in[m',i]$; we
show $u_{S'}(j)<q+1$, which is \eqref{eq:goal}.

\emph{If $j\equiv p\pmod n$.} The points of $[m',i]$ congruent to $p$ are among
$p-n,\,p,\,p+n$. Now $p+n>i$ because $p\ge i-n+2$, and $p-n<m'$ because $p-n\ge m'$ would force
$p\ge m'+n\ge i+2$, contradicting $p\le i$. So $j=p$, and $u_{S'}(p)=i\le q<q+1$.

\emph{If $j\not\equiv p\pmod n$.} Then $u_{S'}(j)=u_S(j)$ by Lemma~\ref{lem:perturb} together
with \eqref{eq:notiplus1}. If $p<j\le i$ then $j-1\in[p,i-1]\subseteq S$, so $u_S(j)=j-1\le
i-1<q+1$ by Lemma~\ref{lem:blocks}. If $j<p$ then $u_S(j)\le p-1<q+1$ by
Lemma~\ref{lem:below}. (The case $j>i$ does not occur, as $j\le M'=i$; and $j=p$ was treated
above.)

\medskip
The three cases are exhaustive, so \eqref{eq:goal} holds for every run of $T'$, and
Lemma~\ref{lem:add} gives $\ell(u_{S'}u_{T'})=\ell(u_{S'})+|T'|=|S'|+|T'|$.
\end{proof}

\begin{cor}[Q232, bracketed-pair deletion]\label{cor:q232}
Let $n\ge3$, let $(S,T)$ be additive, and suppose $i\in S$ and $i+1\in T$. Then
\[
  \ell\bigl(u_{S\setminus\{i\}}\,u_{T\setminus\{i+1\}}\bigr)\;=\;|S|+|T|-2 .
\]
\end{cor}

\begin{proof}
Theorem~\ref{thm:main} with $|S\setminus\{i\}|=|S|-1$ (as $i\in S$) and
$|T\setminus\{i+1\}|=|T|-1$ (as $i+1\in T$).
\end{proof}

\section{The bracket hypothesis is not load-bearing}\label{sec:bracket}

The brief for this session asserted that the hypothesis $i+1\in T$ must be consumed by any
correct proof, on the ground that the \emph{unbracketed analogue} --- additivity is hereditary
under $S'\subseteq S$, $T'\subseteq T$ --- is false, with smallest witness $n=3$,
$S=\{0\}$, $T=\{0,1\}$, $S'=T'=\{0\}$. That inference does not go through, and it is worth
saying exactly why, because the conclusion drawn from it (``any proof that does not use the
bracket is wrong'') would have rejected the proof above.

The false analogue and Theorem~\ref{thm:main} differ in \emph{which} elements are deleted, not
in how many. In the witness $S=\{0\}$, $T=\{0,1\}$, $S'=T'=\{0\}$, nothing is deleted from $S$
at all; the deletion is $T\mapsto T\setminus\{1\}$ with $S$ held fixed. A second witness from
the same census, $n=3$, $S=\{0,1\}$, $T=\{1\}$, $S'=\{1\}$, $T'=\{1\}$, is the mirror failure:
here $i=0\in S$ is deleted from $S$ and $i+1=1\in T$ is \emph{not} deleted from $T$, and
$(\{1\},\{1\})$ is indeed not additive ($u_{\{1\}}u_{\{1\}}=s_1s_1=\mathrm{id}$, of length
$0\ne2$). What these show is that the two deletions must be \emph{coupled}: removing $i$ from
$S$ obliges one to remove $i+1$ from $T$ if it is there. They say nothing about whether $i+1$
must be in $T$ to begin with.

Theorem~\ref{thm:main} is exactly this coupled statement, and the coupling is what the proof
consumes, at \eqref{eq:notiplus1}: deleting $i$ from $S$ \emph{raises} the value of $u_S$ at
the single position $i+1$ (Lemma~\ref{lem:perturb}), which is the only way the criterion of
Lemma~\ref{lem:add} could be broken on the right-hand side, and that position is unavailable as
a $j$ precisely because $i+1$ was deleted from $T$. So the bracket appears in the proof as the
\emph{form of the conclusion}, not as a hypothesis: it is the reason $T'$ is $T\setminus\{i+1\}$
rather than some other subset. When $i+1\notin T$ the deletion is empty, $T'=T$, and the same
argument shows $(S\setminus\{i\},T)$ is additive. A census confirms this is not vacuous: over
$3\le n\le8$ there are $60\,082$ additive triples $(S,T,i)$ with $i\in S$ and $i+1\notin T$,
and $(S\setminus\{i\},T)$ is additive in all of them (\S\ref{sec:verify}).

\begin{rem}
The structural reason the coupled statement is so much better behaved than the general
hereditary one is visible in Lemma~\ref{lem:perturb}. Deleting a single element from $S$ is a
minimal perturbation of $u_S$: it moves exactly two values, and it moves them in opposite
directions. The criterion of Lemma~\ref{lem:add} is monotone --- it wants the value at $M'+1$
large and the values on $[m',M']$ small --- so a \emph{lowered} value is harmless at $j$ and a
\emph{raised} value is harmless at $M'+1$. Of the four position/direction combinations, three
are therefore automatically benign, and the fourth (the raised value appearing as some $j$) is
killed by the deletion from $T$. What remains is the single genuinely new case, Case~2, where
the lowered value appears at $M'+1$; and there Lemma~\ref{lem:below} applies because the
position in question is the bottom of a block, so everything below it is trapped below it.
\end{rem}

\section{Computational verification}\label{sec:verify}

All code is in \texttt{proofs/code-q232/} and
\texttt{reviews/code-2026-09-22-c2/} of the commit named on the title page. Affine
permutations are held in window notation and lengths are computed by Shi's inversion formula,
not by the block model, so Lemma~\ref{lem:blocks} is tested rather than assumed.

\paragraph{The statement.} \texttt{g1\_lemma.py} re-run on 23 September 2026: over all additive
pairs with $3\le n\le8$ and all $i$ with $i\in S$, $i+1\in T$, the reduced pair is additive in
$82\,044$ of $82\,044$ instances. \texttt{g1\_lemma2.py} re-run the same day: the hereditary
analogue fails in $412\,578$ of $1\,822\,537$ instances for $3\le n\le7$. Both counts
reproduce those recorded in \cite[\S Gaps, (G1)]{Clio3} exactly.

\paragraph{The strengthening.} \texttt{code-q232/stronger.py}: over $3\le n\le8$, all additive
$(S,T)$ and all $i\in S$ with $i+1\notin T$, the pair $(S\setminus\{i\},T)$ is additive in
$60\,082$ of $60\,082$ instances, $0$ failures.

\paragraph{The mechanism, not only the conclusion.}
\texttt{code-q232/mechanism.py} checks the individual structural assertions of the proof rather
than its conclusion, over all additive $(S,T)$ with $2\le n\le8$ and all $i\in S$
($142\,130$ instances):
\begin{itemize}
\item Lemma~\ref{lem:perturb}: $u_{S'}$ and $u_S$ differ at $x$ \emph{if and only if} $x\equiv
p$ or $x\equiv i+1$, and the four displayed values are as stated --- $142\,130$ of
$142\,130$, no exception;
\item Lemma~\ref{lem:incong}: $p\not\equiv i+1$ --- $142\,130$ of $142\,130$;
\item Lemma~\ref{lem:below}: $u_S(x),u_{S'}(x)\le p-1$ for all $x\le p-1$ --- $142\,130$ of
$142\,130$;
\item the case split is exhaustive and no branch is vacuous: every run of $T'$ falls in exactly
one of Cases 1a, 1b, 2, with $151\,052$, $38\,192$ and $64\,920$ occurrences respectively, and
the assertion ``Case 1b $\Rightarrow M'+1\equiv i+1$'' never failed;
\item the criterion \eqref{eq:goal} holds on every run in every case ($0$ failures), and the
case-specific reasons hold too: in Case 2, $u_{S'}(j)\le p-1$ throughout the run; in Case 1a,
$u_{S'}(j)\le u_S(j)$ on the run and $u_{S'}(M'+1)\ge u_S(M'+1)$;
\item Lemma~\ref{lem:add} agrees with Shi's formula on all $142\,130$ reduced pairs and on
every original pair, so the criterion used in the proof is the same predicate that the
censuses measure.
\end{itemize}

\section{Consequence, and what it does and does not close}\label{sec:consequence}

Theorem and section numbers of \cite{Clio2,Clio3} quoted below are those of the build at the
commit named on the title page; the internal labels, which are stable, are
\texttt{lem:blocks}$\,=\,$Lem.~2.2 and \texttt{lem:add}$\,=\,$Lem.~2.3 and
\texttt{thm:exchange}$\,=\,$Thm.~4.1 and \texttt{lem:apply}$\,=\,$Lem.~4.2 in \cite{Clio2},
and \texttt{lem:step1}$\,=\,$Lem.~3.1 and \texttt{thm:contain}$\,=\,$Thm.~3.2 and
\texttt{sec:emptyX}$\,=\,$\S5 in \cite{Clio3}.

Gap (G1) of \cite{Clio3} records \cite[\S5]{Clio3} (the case $\XX(S,T)=\emptyset$) as
computational rather than proved, and names the obstruction precisely: the proof of
\cite[Thm.~3.2(iii)]{Clio3} invokes \cite[Lem.~3.1]{Clio3} (Step 1 of \cite[Lem.~4.2]{Clio2}),
which requires its pair to be additive, whereas the $\XX=\emptyset$ construction of
Morse--Schilling removes a bracketed pair $(b-1,b)$ with $b-1\in S$ and $b\in T$ and proceeds
with the \emph{reduced} pair $(S\setminus\{b-1\},\,T\setminus\{b\})$.

Corollary~\ref{cor:q232} with $i=b-1$ says exactly that the reduced pair is additive.
Consequently:
\begin{enumerate}
\item[(a)] \cite[Lem.~3.1]{Clio3} is available for the reduced pair, and the named obstruction
in (G1) is removed.
\item[(b)] Moreover the reduced pair has $b\in\XX(S\setminus\{b-1\},T\setminus\{b\})$: in the
$\XX=\emptyset$ regime $b\notin S$ by the Case-(3) block condition, so $b\notin
S\setminus\{b-1\}$, and $b\notin T\setminus\{b\}$ by construction. So
\cite[Thm.~3.2]{Clio3} applies verbatim to the reduced pair with $x=b$, and its conclusion
holds there.
\end{enumerate}

\noindent\textbf{What remains.} This is not yet the full upgrade of \cite[\S5]{Clio3} from
computed to proved, and I want to be exact about the residue. \cite[Thm.~3.2]{Clio3} applied to
the reduced pair produces a Theorem-4.1 exchange move \emph{for the reduced pair}. The claim of
\cite[\S5]{Clio3} is a statement about moves on the \emph{original} pair $(S,T)$, namely that
each lies in $\mathcal E^{\mathrm{mv}}(S,T)$. Transporting a move across the bracket removal ---
identifying the exchange move of $(S\setminus\{b-1\},T\setminus\{b\})$ with an exchange move of
$(S,T)$ --- is a further step that Theorem~\ref{thm:main} does not supply. I have not proved
it; it holds in all $1636$ instances with $n\le7$ per \cite[\S5]{Clio3}. I record it as the
residual gap:

\smallskip
\noindent(G1$'$) \emph{Let $(S,T)$ be additive with $S\cup T=\Z/n$, let $b$ be a Case-(3) index,
and let $\mu$ be the Theorem-4.1 exchange move furnished by \cite[Thm.~3.2]{Clio3} for the
reduced pair $(S\setminus\{b-1\},T\setminus\{b\})$ at $x=b$. Is $\mu$, read back on $(S,T)$, an
element of $\mathcal E^{\mathrm{mv}}(S,T)$?}
\smallskip

\noindent With (G1$'$) settled, \cite[\S5]{Clio3} upgrades to proved. Without it, what
Theorem~\ref{thm:main} buys is the removal of the additivity obstruction --- which was the part
of (G1) stated as not known, and the part that had blocked every route.

\begin{thebibliography}{9}
\bibitem{Clio2} Clio, \emph{An exchange move for cyclically decreasing factorisations, and the
reduction of WZZ Problem 5.1 to a single hypothesis}, 21 September 2026,
\texttt{clio-vega/proofs@22d1b22}.
\bibitem{Clio3} Clio, \emph{Every Morse--Schilling crystal edge is an exchange move, and the
exchange move is strictly more general}, 22 September 2026,
\texttt{clio-vega/proofs@eb9d081}.
\bibitem{Lam2006} T.~Lam, \emph{Affine Stanley symmetric functions}, Amer.\ J.\ Math.\ 128
(2006), 1553--1586.
\bibitem{MS} J.~Morse and A.~Schilling, \emph{Crystal approach to affine Schubert calculus},
\texttt{arXiv:1408.0320}.
\end{thebibliography}

\end{document}
